\documentclass[11pt]{article}

\usepackage{amsmath}
\usepackage{amssymb}
\usepackage{bm}
\usepackage{geometry}
\usepackage{graphicx}
\usepackage{hyperref}

\geometry{margin=1in}

\title{An Extended Kalman Filter for Skid-Steer Robot Odometry with Online Effective Wheel Separation Estimation}

\author{
Stuart G. Johnson \and ChatGPT (OpenAI)
}

\date{\today}

\begin{document}

\maketitle

\begin{abstract}
This note derives an Extended Kalman Filter (EKF) for odometry estimation
of a skid-steer robot modeled approximately as a differential-drive robot.
The filter estimates planar pose, gyroscope bias, and an effective wheel
separation. The effective wheel separation is not interpreted as a purely
geometric quantity; rather, it is a compact parameter that absorbs some
of the mismatch between ideal differential-drive kinematics and the
actual skid-steer motion.
\end{abstract}

\section{Motivation}

A four-wheel skid-steer robot is often modeled as a differential-drive
robot. This approximation is convenient, but it is imperfect. During
turns, the wheels scrub against the ground. This violates the ideal
nonholonomic differential-drive model.

The goal of this filter is to provide proprioceptive odometry suitable for ROS2
navigation. The odometry will drift over time, but it should
be locally smooth and track true robot motion. The filter uses wheel
velocity estimates and gyroscope measurements. It also estimates an effective
wheel separation in order to adapt the differential-drive model to the
observed robot behavior.

\section{State Vector}

The EKF state at time $k$ is

\begin{equation}
\mathbf{x}_k =
\begin{bmatrix}
x_k \\
y_k \\
\theta_k \\
b_{g,k} \\
b_k
\end{bmatrix}.
\end{equation}

The state components are:

\begin{itemize}
\item $x_k,y_k$: robot position in the odometry frame,
\item $\theta_k$: robot heading,
\item $b_{g,k}$: gyroscope yaw-rate bias,
\item $b_k$: effective wheel separation.
\end{itemize}

The effective wheel separation $b_k$ should not be interpreted too
literally as the physical distance between wheels. For a skid-steer
robot, it is an effective kinematic parameter that can absorb wheel slip,
tire compliance, load transfer, and other unmodeled effects.

\section{Inputs and Measured Quantities}

Let the measured left and right wheel linear velocities be

\begin{equation}
\mathbf{u}_k =
\begin{bmatrix}
v_{L,k} \\
v_{R,k}
\end{bmatrix}.
\end{equation}

The differential-drive forward velocity estimate is

\begin{equation}
v_k =
\frac{v_{R,k}+v_{L,k}}{2}.
\end{equation}

The wheel-derived yaw-rate estimate is

\begin{equation}
\omega_k =
\frac{v_{R,k}-v_{L,k}}{b_k}.
\end{equation}

The gyroscope measurement is denoted

\begin{equation}
z_{g,k}.
\end{equation}

For this EKF, the wheel and gyroscope data are averaged over
common robot-side sampling intervals and then emitted to the ROS2 system.
Thus, the filter should process synchronized wheel and IMU packets as
one discrete update step.

\section{Continuous-Time Process Model}

The continuous-time kinematic model is

\begin{align}
\dot{x} &= v \cos\theta, \\
\dot{y} &= v \sin\theta, \\
\dot{\theta} &= \frac{v_R-v_L}{b}, \\
\dot{b}_g &= 0, \\
\dot{b} &= 0.
\end{align}


\section{Discretization}

The Euler update process model is

\begin{equation}
\mathbf{x}_{k+1} = f(\mathbf{x}_k, \mathbf{u}_k) + \textbf{w}_k
\end{equation}

where $\textbf{w}_k$ is the process noise. $f$ is defined as:

\begin{align}
\mathbf{x}_{k+1}&= f(\mathbf{x}_k, \mathbf{u}_k)
\\
x_{k+1}
&=
x_k
+
\Delta t_k
v_k
\cos\theta_k,
\\
y_{k+1}
&=
y_k
+
\Delta t_k
v_k
\sin\theta_k,
\\
\theta_{k+1}
&=
\theta_k
+
\Delta t_k
\omega_k,
\\
b_{g,k+1}
&=
b_{g,k},
\\
b_{k+1}
&=
b_k.
\end{align}

\section{Process Jacobian}

The EKF process Jacobian is

\begin{equation}
F_k =
\left.
\frac{\partial f}{\partial \mathbf{x}}
\right|_{\hat{\mathbf{x}}_{k|k},\mathbf{u}_k}.
\end{equation}

The full state Jacobian is

\begin{equation}
F_k =
\begin{bmatrix}
1 & 0 & -\Delta t_k v_k sin\theta_k & 0 & 0 \\
0 & 1 & \Delta t_k v_k cos\theta_k & 0 & 0 \\
0 & 0 & 1 & 0 &
-\Delta t_k \frac{\omega_k}{b_k}
\\
0 & 0 & 0 & 1 & 0
\\
0 & 0 & 0 & 0 & 1
\end{bmatrix}.
\end{equation}


\section{Input and Process Noise Propagation}

Let the wheel speed covariance be

\begin{equation}
\Sigma_u =
\begin{bmatrix}
\sigma_{v_L}^2 & \sigma_{v_L v_R} \\
\sigma_{v_L v_R} & \sigma_{v_R}^2
\end{bmatrix}.
\end{equation}

If the left and right wheel speed errors are assumed
uncorrelated, then

\begin{equation}
\Sigma_u =
\begin{bmatrix}
\sigma_{v_L}^2 & 0 \\
0 & \sigma_{v_R}^2
\end{bmatrix}.
\end{equation}

The input Jacobian is

\begin{equation}
G_k
=
\left.
\frac{\partial f}{\partial \mathbf{u}}
\right|_{\hat{\mathbf{x}}_{k|k},\mathbf{u}_k}.
\end{equation}

which evaluates to:

\begin{equation}
G_k
=
\begin{bmatrix}
\Delta t_k\frac{1}{2}cos\theta_k &
\Delta t_k\frac{1}{2}cos\theta_k)
\\
\Delta t_k\frac{1}{2}sin\theta_k  &
\Delta t_k\frac{1}{2}sin\theta_k 
\\
-\frac{\Delta t_k}{b_k} &
\frac{\Delta t_k}{b_k}
\\
0 & 0
\\
0 & 0
\end{bmatrix}.
\end{equation}

Random-walk process noise for gyroscope bias and effective wheel separation is

\begin{equation}
Q_k
=
\begin{bmatrix}
0 & 0 & 0 & 0 & 0 \\
0 & 0 & 0 & 0 & 0 \\
0 & 0 & 0 & 0 & 0 \\
0 & 0 & 0 & \sigma_{bg}^2 \Delta t_k & 0 \\
0 & 0 & 0 & 0 & \sigma_b^2 \Delta t_k
\end{bmatrix}.
\end{equation}

Thus the total process covariance contribution is

\begin{equation}
Q_{k}
=
G_k
\Sigma_u
G_k^{T}
+
Q_k
\end{equation}

\section{Gyroscope Measurement Model}

The gyroscope measurements return yaw rate plus bias:

\begin{equation}
z_{g,k}
=
\omega_k
+
b_{g,k}
+
\nu_k,
\end{equation}

where $\nu_k$ is zero-mean gyroscope measurement noise.

The measurement function is

\begin{equation}
h(\mathbf{x}_k,\mathbf{u}_k)
=
\frac{v_{R,k}-v_{L,k}}{b_k}
+
b_{g,k}.
\end{equation}

The measurement Jacobian is

\begin{equation}
H_k
=
\left.
\frac{\partial h}{\partial \mathbf{x}}
\right|_{\hat{\mathbf{x}}_{k+1|k},\mathbf{u}_k}.
\end{equation}

Explicitly,

\begin{equation}
H_k
=
\begin{bmatrix}
0 &
0 &
0 &
1 &
-\frac{v_{R,k}-v_{L,k}}{b_k^2}
\end{bmatrix}.
\end{equation}

The gyroscope measurement covariance is

\begin{equation}
R =
\sigma_g^2.
\end{equation}

\section{EKF Algorithm}

At each synchronized measurement, the EKF performs the following
operations:

\subsection{Prediction}

Given the posterior estimate from the previous step,

\begin{equation}
\hat{\mathbf{x}}_{k|k},
\end{equation}

and covariance

\begin{equation}
P_{k|k},
\end{equation}

compute

\begin{equation}
\hat{\mathbf{x}}_{k+1|k}
=
f(\hat{\mathbf{x}}_{k|k},\mathbf{u}_k).
\end{equation}

Compute $F_k$ and $G_{u,k}$ at
$\hat{\mathbf{x}}_{k|k}$ and $\mathbf{u}_k$. Then propagate covariance:

\begin{equation}
P_{k+1|k}
=
F_k
P_{k|k}
F_k^T
+
Q_k.
\end{equation}

\subsection{Update}

The gyroscope innovation is

\begin{equation}
r_k
=
z_k
-
h(\hat{\mathbf{x}}_{k+1|k},\mathbf{u}_k).
\end{equation}
The innovation covariance is
\begin{equation}
S_k
=
H_k
P_{k+1|k}
H_k^{T}
+
R.
\end{equation}
The Kalman gain is
\begin{equation}
K_{g,k}
=
P_{k+1|k}
H_{g,k}^{T}
S_{g,k}^{-1}.
\end{equation}
The posterior state estimate is
\begin{equation}
\hat{\mathbf{x}}_{k+1|k+1}
=
\hat{\mathbf{x}}_{k+1|k}
+
K_{g,k}r_{g,k}.
\end{equation}
The covariance update is:
\begin{equation}
P_{k+1|k+1}
=
(I-K_kH_k)
P_{k+1|k}
\end{equation}

\section{ROS2 Odometry Covariances}

The ROS2 \texttt{nav\_msgs/Odometry} message contains pose and twist
covariances.

The pose is expressed in the frame named by \texttt{header.frame\_id},
typically \texttt{odom}. The twist is expressed in the frame named by
\texttt{child\_frame\_id}, typically \texttt{base\_link}.

\subsection{Pose Covariance}

The planar pose covariance is obtained directly from the EKF covariance
matrix. The entries corresponding to $x$, $y$, and yaw are inserted into
the ROS2 $6\times 6$ pose covariance matrix with ordering

\begin{equation}
[x,\ y,\ z,\ roll,\ pitch,\ yaw].
\end{equation}

For example,

\begin{align}
\Sigma_{pose}(x,x) &= P(x,x), \\
\Sigma_{pose}(y,y) &= P(y,y), \\
\Sigma_{pose}(yaw,yaw) &= P(\theta,\theta), \\
\Sigma_{pose}(x,yaw) &= P(x,\theta), \\
\Sigma_{pose}(y,yaw) &= P(y,\theta).
\end{align}

\subsection{Twist Covariance}

The state definition does not include the twist. Therefore twist covariance must be derived from the uncertainty in the reported twist estimate.

The twist estimate returned to ROS2 is:

\begin{align}
v &= \frac{v_L+v_R}{2}, \\
\omega &= \frac{v_R-v_L}{b},
\end{align}

The ROS2 twist covariance uses ordering

\begin{equation}
[v_x,\ v_y,\ v_z,\ \omega_x,\ \omega_y,\ \omega_z].
\end{equation}

Twist covariance for the non-zero components $(v_x,\omega_z)$ is obtained by propagation of errors:

\begin{equation}
\Sigma_{\mathrm{twist}} (v_x,\omega_z)
\approx
J
\begin{bmatrix}
\sigma_{v_L}^2 & 0 & 0 \\
0 & \sigma_{v_R}^2 & 0 \\
0 & 0 & \sigma_b^2
\end{bmatrix}
J^T
\end{equation}
where:
\begin{equation}
J =
\begin{bmatrix}
\dfrac{\partial v}{\partial v_L} &
\dfrac{\partial v}{\partial v_R} &
\dfrac{\partial v}{\partial b} \\[8pt]

\dfrac{\partial \omega}{\partial v_L} &
\dfrac{\partial \omega}{\partial v_R} &
\dfrac{\partial \omega}{\partial b}
\end{bmatrix}
=
\begin{bmatrix}
\dfrac{1}{2} & \dfrac{1}{2} & 0 \\[10pt]

-\dfrac{1}{b} & \dfrac{1}{b} &
-\dfrac{(v_R-v_L)}{b^2}
\end{bmatrix}
\end{equation}
Note that $\Sigma_{\mathrm{twist}} (v_x,\omega_z)$ is a 2x2 matrix which must be distributed into the twist covariance.

ROS2 does not generally like covariance diagonals which are zero. For twist and state covariance, otherwise zero diagonals are set to $9999.0$.


\end{document}
